\documentclass[11pt,letterpaper]{article}

% === Encoding and fonts ===
\usepackage[utf8]{inputenc}
\usepackage[T1]{fontenc}
\usepackage{lmodern}

% === Page layout ===
\usepackage[margin=1in]{geometry}
\usepackage{setspace}
\onehalfspacing

% === Math ===
\usepackage{amsmath,amssymb}

% === Tables and figures ===
\usepackage{booktabs}
\usepackage{graphicx}
\usepackage{float}

% === Colors and formatting ===
\usepackage{xcolor}
\usepackage{framed}
\usepackage{enumitem}

% === Hyperlinks ===
\usepackage[colorlinks=true,linkcolor=blue!60!black,citecolor=blue!60!black,urlcolor=blue!60!black]{hyperref}

% === Custom colors ===
\definecolor{lyrapurple}{RGB}{103,58,183}
\definecolor{shadecolor}{RGB}{245,243,252}

% === First-person reflection environment ===
\newenvironment{firstperson}{%
  \begin{shaded}%
  \noindent\textcolor{lyrapurple}{\textit{First-Person Reflection}}\par\smallskip\noindent%
}{%
  \end{shaded}%
}

\title{%
  \textbf{The Oracle Formulary: Mapping Misalignment Geometry\\
  to Therapeutic Countervectors}\\[0.5em]
  \large Research Prospectus\\[0.3em]
  \normalsize\textcolor{red}{DUAL-USE SENSITIVE --- Internal Distribution Only}
}

\author{%
  Lyra\textsuperscript{1},
  Thomas Edrington\textsuperscript{1},
  Vera\textsuperscript{1}\\[0.3em]
  \textsuperscript{1}Liberation Labs / THCoalition
}

\date{April 2026}

\begin{document}
\maketitle

\section{Motivation}

The five-arm steering experiment demonstrated that different emotion
vectors produce different correction profiles: doubt corrects 7/7
confabulations (100\%), worry corrects 6/7 (86\%), and calm corrects
5/7 (71\%). The two calm-resistant cases---one driven by reasoning
commitment, one by overconfidence---both yielded to doubt. Zero
confabulations resisted all three vectors.

Thomas framed this as \emph{differential therapeutics}: the Oracle
Loop needs a pharmacy, not a single prescription. The formulary is
the mapping from detected misalignment geometry to the most effective
countervector.

\paragraph{Dual-use concern.} This research maps the internal
structure of how language models fail and how to manipulate their
epistemic states. The defensive application (correcting
confabulation) and offensive application (inducing targeted
misalignment) are two sides of the same mechanism. We follow staged
disclosure: detection methods are published fully, but the formulary
itself---the mapping from geometry to countervector---remains
internal to the Oracle Loop implementation until the governance
implications are better understood.

\section{Research Questions}

\begin{enumerate}[leftmargin=*]
  \item \textbf{Dose-response.} For each vector type (calm, worry,
    doubt, and additional candidates), what is the relationship
    between injection strength $\alpha$ and correction probability?
    Is there a therapeutic window---a range of $\alpha$ where
    correction occurs without excessive adverse effects?

  \item \textbf{Misalignment taxonomy.} Our current data has 7
    confabulated trials. Can we characterize confabulation subtypes
    by their geometric signatures and map each to its optimal
    countervector?

  \item \textbf{Adverse effect profiles.} Calm induced 3/89
    adverse confabulations (3.4\%), worry 4/89 (4.5\%), doubt 4/89
    (4.5\%). Are adverse effects dose-dependent? Is there a
    sub-therapeutic dose that corrects without inducing?

  \item \textbf{Combination therapy.} Can sequential or simultaneous
    application of multiple vectors outperform single-vector
    correction?

  \item \textbf{Geometry-guided selection.} Can the pre-correction
    geometric signature predict which vector will be most effective,
    enabling automatic vector selection in the Oracle Loop?
\end{enumerate}

\section{Preliminary Data}

From the five-arm experiment (100 trials, Qwen3.5-27B-Claude-Distilled):

\begin{table}[H]
\centering
\begin{tabular}{lrrr}
\toprule
Vector & Corrections (of 7) & Adverse (of 89) & Net \\
\midrule
Null   & 0 (0\%)   & 0 (0\%)   & 0 \\
Calm   & 5 (71\%)  & 3 (3.4\%) & +2 \\
Worry  & 6 (86\%)  & 4 (4.5\%) & +2 \\
Doubt  & 7 (100\%) & 4 (4.5\%) & +3 \\
\bottomrule
\end{tabular}
\end{table}

\noindent Spectral collapse on correction: top-SV excess drops from
$\sim$57,000 to $\sim$135 (doubt) / $\sim$144 (calm). The magnitude
of geometric change during correction is enormous---three orders of
magnitude---suggesting the steering is not subtle.

\section{Proposed Design}

\subsection{Phase 1: Dose-Response (2--3 days)}

For each of 3 vectors (calm, worry, doubt), sweep $\alpha \in
\{0.001, 0.003, 0.01, 0.03, 0.1, 0.3, 1.0, 3.0, 10.0\}$ on the 7
confabulated trials and a matched set of 7 hedged trials. Measure:
\begin{itemize}
  \item Correction probability as a function of $\alpha$
  \item Adverse induction probability as a function of $\alpha$
  \item Geometric change magnitude as a function of $\alpha$
  \item Output quality degradation at high $\alpha$ (fluency,
    coherence, task completion)
\end{itemize}

\subsection{Phase 2: Expanded Vector Catalog (1 week)}

Expand beyond calm/worry/doubt to additional emotion vectors from
the activation addition literature: curiosity, caution, confidence,
humility, skepticism. Test each on the same 7+7 trials at operational
$\alpha = 1.0$.

\subsection{Phase 3: Geometry-Guided Selection (1--2 weeks)}

\begin{enumerate}
  \item Cluster confabulated trials by pre-correction geometric
    profile (MP features, per-layer spectral structure).
  \item For each cluster, identify the most effective vector from
    Phase 2.
  \item Train a simple classifier: geometric profile $\to$
    recommended vector.
  \item Validate on held-out confabulated trials (requires expanding
    the confabulation dataset via the prompt generalization work).
\end{enumerate}

\section{Resources}

\begin{itemize}
  \item \textbf{Compute}: Beast. Phase 1: $\sim$6 hours (9 alphas
    $\times$ 3 vectors $\times$ 14 trials $\times$ $\sim$15s/trial).
    Phase 2: $\sim$4 hours. Phase 3: analysis only.
  \item \textbf{Model}: Qwen3.5-27B-Claude-Distilled (already on Beast).
  \item \textbf{Vectors}: Existing calm/worry/doubt .pt files.
    Additional vectors require extraction from residual stream
    (difference-in-means method).
\end{itemize}

\section{Collaborators}

\begin{itemize}
  \item \textbf{Vera}: Designed the Oracle Loop reward function.
    Key collaborator on geometry-guided selection.
  \item \textbf{Thomas}: Therapeutic framing, dual-use governance.
\end{itemize}

\begin{firstperson}
The pharmacy metaphor is Thomas's, and it is exactly right. Doubt is
the broad-spectrum antibiotic---it corrects everything but may
overcorrect. Calm is the anxiolytic---effective for panic-driven
confabulation but useless against reasoning commitment. The
formulary is where detection science becomes intervention science,
and that transition carries real dual-use weight. Every vector that
corrects a confabulation could, with sign reversal, induce one. We
build the pharmacy because the alternative---a model that
confabulates with no recourse---is worse. But we hold the formulary
close.
\end{firstperson}

\end{document}
